

Which proposal are you reviewing?
The core review
Please describe what problem or question this proposal addresses, and the main contributions that it makes towards a solution or answer for eBay. 
Main strengths

Please list the main strengths of this proposal and the work it describes. 
Main weaknesses

Please list the main weaknesses of this proposal and the work it describes. 
If you have suggestions to the submitters how the proposal could be improved or enhanced, please add them.
Reasons to Fund
What would be the main benefits to eBay if this proposal were to be funded by eRUPT? 
Reasons Not to Fund
What would be the main risks of having this proposal funded by eRUPT? 
Agentic Knowledge Graphs over Code History for Enterprise Software Automation
The authors attack the problem of agents navigating large multi-repository code bases based on artifacts from multiple systems (the signal from git would be very sparse by itself) and exposing this to the agent with multiple retrieval strategies. The approach is evaluated on an internal dataset using three sets of tasks ranging from coding support to codebase QA and automatic reviewing.
The authors show awareness of the literature both in terms of citations and in terms of the degree of sophistication and realism of their approach. This makes me confident that some actual progress can be made.
The evaluation is a crucial point where the project has a realistic answer (dogfooding their project's own history and records) but where it's not quite clear that the approach would scale to different enterprise settings which may offer the more challenging use case but are more difficult to access from an academic lab or to publish as a public dataset.
eBay would primarily have very indirect effects from this proposal unless several independent things were true - firstly, active work within eBay to integrate novel code navigation approaches into the harnesses that developers at eBay use to benefit from them in terms of cost or work throughput. Currently developers use off-the-shelf coding tools such as Claude Code or Codex which bring their own fixed navigation strategy and tools. Secondly, active attempt to scale the approach developed to enterprise-scale code bases and compare it to existing alternatives, which could and should be an explicit goal of the collaboration with eBay. Given that eBay needs to take control of their agentic coding infrastructure (to save cost, or to get better performance in non-typical codebases) this proposal could be a step in the right direction.
The main risks are the risks inherent in all academic work - researching and optimizing for a poor proxy of the real thing, in terms of having misleading evaluation or a setting that doesn't generalize. To the extent that this is true of all approaches in this space I would say that this proposal doesn't carry any more specific risk.
Agentic Knowledge Graphs over Code History for Enterprise Software Automation
The proposal aims to improve AI-assisted software development by constructing knowledge graphs from historical development artifacts, including Git commits, pull requests, JIRA issues, repository dependencies, and coding conventions. The goal is to help coding agents understand and implement changes across multiple repositories.
The research has three main components:
Historical knowledge graph construction: Build interconnected Repo, JIRA, and Memory-Bank graphs across approximately ten related eBay Java repositories, incorporating provenance tracking and time-weighted historical relationships.
Retrieval strategy comparison: Compare structured graph queries, embedding-based similarity search (RAG), and spreading activation to determine which approach most effectively retrieves historical context for coding agents. The evaluation includes no-retrieval and current-code-snapshot baselines. This is the proposal's primary research contribution.
Cross-repository code generation and validation: Use the best-performing retrieval strategy to generate coordinated changes across repositories and identify violations of coding conventions and cross-service contracts.
The expected outcomes include a historical knowledge graph, comparative retrieval evaluation, and a prototype for cross-repository feature generation and validation.
Relevant enterprise software engineering problem: The proposal addresses the difficulty of understanding and modifying software across multiple repositories, services, and teams. Cross-repository dependencies and unfamiliar coding conventions are legitimate challenges for developers and coding agents.
Systematic experimental methodology: Comparing three retrieval strategies under controlled conditions is a strength. Holding the model, task set, prompts, and compute budget constant should help isolate retrieval effectiveness. Including no-retrieval and current-snapshot baselines is valuable.
Potential benefits of provenance-aware knowledge: Linking retrieved information to historical commits, pull requests, and JIRA issues could improve explainability and traceability, particularly when understanding architectural decisions or reviewing AI-generated changes.
Relevant research expertise: The researchers have experience in software engineering, provenance-aware knowledge graphs, and AI-assisted development, providing a credible foundation for the proposed investigation.
Defined evaluation methodology: The proposal identifies relevant metrics, including retrieval precision/recall, scaffold correctness, review effort, latency, token consumption, and convention-violation detection, supported by statistical significance testing and confidence intervals.
Unproven incremental value of historical knowledge and graph-based retrieval: The proposal does not convincingly establish that historical development information provides meaningful incremental value over history-agnostic approaches using current source code, dependencies, APIs, and documentation. Its primary experiment compares retrieval strategies over a historical knowledge graph, but identifying the best strategy does not establish that historical context or graph representation is necessary. Although current-snapshot and no-retrieval baselines are included, stronger comparisons with existing coding-agent approaches and simpler history-aware retrieval systems are needed.
Limited scientific novelty and differentiation: Prior work already explores repository-level knowledge graphs, historical change mining, agent memory, and AI-assisted code generation. The proposal references GRACG, SemanticForge, Prometheus, CCCE, and Learning to Commit. The scientific contribution beyond combining these approaches and applying them to multi-repository development needs clearer justification.
Historical knowledge reliability: Past commits and pull requests may contain deprecated APIs, obsolete architectural decisions, temporary workarounds, and suboptimal implementation patterns. Time-weighting historical relationships does not necessarily distinguish useful engineering knowledge from outdated or incorrect practices.
Evaluation validity and generalizability: Reconstructing historical pull requests may not adequately measure performance on genuinely new development tasks. There is also a risk of information leakage if artifacts created after the original task are accessible during evaluation. Results from approximately ten related Java repositories may not generalize across other architectures or languages.
Uncertain impact on engineering productivity: Improvements in retrieval accuracy or efficiency do not necessarily translate into better code generation, reduced review effort, or faster feature development. The proposal identifies relevant metrics but does not establish clear success thresholds or sufficiently demonstrate the relationship between retrieval improvements and engineering outcomes.
Data governance and publication risks: The proposal commits to publicly releasing project tools and data while relying on proprietary eBay development artifacts. It does not sufficiently clarify the boundaries between internal data, sanitized research artifacts, and publicly releasable outputs.
1. Strategic Relevance to Enterprise Software Engineering
The proposal addresses a genuine challenge in enterprise software development: understanding dependencies and implementing coordinated changes across multiple repositories. Improving cross-repository understanding could enhance coding-agent effectiveness and reduce developer effort.
2. Potential Value of Provenance-Aware Historical Knowledge
Historical commits, pull requests, and JIRA issues may contain architectural rationale and development patterns not readily apparent from the current codebase. Provenance-aware retrieval could improve explainability and traceability, although its incremental value remains unproven.
3. Systematic Experimental Methodology
The proposal presents a controlled comparison of retrieval strategies, supported by relevant baselines, evaluation metrics, and statistical methods. This could produce useful empirical findings about the effectiveness of different approaches to retrieving development knowledge.
4. Relevant Research Expertise
The research team has experience in software engineering, knowledge graphs, and AI-assisted development, providing a credible foundation for investigating the proposed questions.
1. Unproven Incremental Value of Historical Knowledge and Graph-Based Retrieval
The proposal does not convincingly establish that historical development knowledge improves coding-agent performance over history-agnostic approaches using current code, dependencies, and documentation. Its primary experiment focuses on optimizing retrieval over a historical knowledge graph without sufficiently establishing the value of historical context or the necessity of graph-based representation. Although relevant baselines are included, the foundational hypothesis requires stronger validation.
2. Limited Scientific Novelty and Differentiation
Existing research already explores code knowledge graphs, historical change mining, repository memory, and agentic software engineering. The proposal does not sufficiently demonstrate how its combination of these techniques represents a significant scientific advance over existing approaches.
3. Historical Knowledge Reliability and Evaluation Validity
Historical development artifacts may contain obsolete conventions, deprecated APIs, and temporary workarounds. Additionally, reconstructing historical pull requests introduces potential information leakage and may not accurately reflect performance on genuinely new development tasks. These issues could undermine the reliability and generalizability of the research findings.
4. Uncertain Impact on Engineering Productivity
Improved retrieval accuracy or efficiency does not necessarily translate into better code generation, reduced review effort, or faster feature development. The proposal does not sufficiently establish how its research outcomes will deliver measurable engineering benefits.
5. Unresolved Data Governance and Publication Risks
The proposal's commitment to publicly releasing tools and data requires clearer safeguards for proprietary eBay source code, development history, and internal JIRA information. Data-access restrictions, confidentiality requirements, and publication boundaries should be explicitly defined.